\section{Affine structure and centered coefficients}
\label{p1:alg:section}

The arithmetic estimates naturally count relations, whereas the
approximation problem concerns a signed difference of laws.
We introduce the exact identities linking these two objects.
They will be used with a stopping rule, so that large connected
families outside the required nullity range never need to be
summed in absolute value.

Throughout this section $I\subseteq\mathcal I_N$, and
$\rho_{\mathrm{rel}}(I)$ is the effective relative nullity.
Write $b_x=\delta_0+p_x\Delta$. The identities below hold for
heterogeneous marginals. On the regular set $G_N$, one has $p_x=p$,
$d_x=L$, and $b_x=b$; their homogeneous specializations are used in
the first reconstruction in the companion.

\subsection{An exact expansion around independent Bernoulli variables}

Set
\[
 Y_x=J_x-p_x,\qquad c(I)=\E\prod_{x\in I}Y_x,
 \qquad c(\varnothing)=1.
\]

\begin{proposition}[Centered Bernoulli expansion]
\label{p1:alg:centred}
For every \(N\),
\begin{equation}
 \Law(Z_N)=\sum_{I\subseteq\mathcal I_N}
       c(I)\Delta^{*|I|}*\mathop{*}_{x\notin I}b_x .
 \label{p1:alg:centred-identity}
\end{equation}
The complementary convolution products are deterministic probabilities.
The terms of sizes zero and two are
\[
 \mathop{*}_{x}b_x+
 \sum_{|e|=2}c(e)\Delta^{*2}*\mathop{*}_{x\notin e}b_x.
\]
Terms of size one vanish.
\end{proposition}

\begin{proof}
For an indicator,
\(\delta_{J_x}=b_x+Y_x\Delta\). Expand the finite convolution
product over \(\mathcal I_N\) and take expectations. No independence
assumption is used in this identity.
\end{proof}

Each \(s\in S_x\) is a square-class vector. Let
\(\chi_s\) be its character on the prime signs, and let
\(\eta_x(s)\in\{\pm1\}\) be the affine sign prescribed by
the run-start event. The local prescription is consistent on
\(\mathcal I_N\), so its affine character \(\eta_x\) is well defined
on the effective space. Thus
\[
 J_x=p_x\sum_{s\in S_x}\eta_x(s)\chi_s,\qquad \eta_x(0)=1.
\]

\begin{lemma}[Full-support Fourier formula]
\label{p1:alg:fourier}
For nonempty \(I\subseteq\mathcal I_N\),
\begin{equation}
 c(I)=\Bigl(\prod_{x\in I}p_x\Bigr)
   \sum_{\substack{s_x\in S_x\setminus\{0\}\ (x\in I)\\
                         \sum_xs_x=0}}
                 \prod_{x\in I}\eta_x(s_x).
 \label{p1:alg:fourier-identity}
\end{equation}
Consequently,
\[
 |c(I)|\le\Bigl(\prod_{x\in I}p_x\Bigr)(2^{\rho_{\mathrm{rel}}(I)}-1)
           \le\Bigl(\prod_{x\in I}p_x\Bigr)2^{\rho_{\mathrm{rel}}(I)}.
\]
If \(c(I)\ne0\), there is a relative relation meeting
every block of \(I\).
\end{lemma}

\begin{proof}
Subtracting \(p_x\) removes the zero character from each
factor. The expectation of a character product is zero
unless the sum of its indices vanishes. The kernel of
\(\bigoplus_{x\in I}S_x\to\sum_{x\in I}S_x\) has
\(2^{\rho_{\mathrm{rel}}(I)}\) elements. The displayed sum omits its
zero element and possibly others. Every affine
coefficient has modulus one.
\end{proof}

\subsection{The nullity that survives affine compatibility}

Define
\begin{equation}
 \rho_*(I)=\max\{\rho_{\mathrm{rel}}(A):A\subseteq I,\ q(A)>0\},
 \qquad q(\varnothing)=1,\quad\rho_{\mathrm{rel}}(\varnothing)=0 .
 \label{p1:alg:compatible-nullity}
\end{equation}
This parameter records the largest rank loss in a
compatible subfamily. It can be smaller than the
nullity of the whole family.

\begin{lemma}[Compatible-nullity majorant]
\label{p1:alg:compatible}
For every \(I\subseteq\mathcal I_N\),
\begin{equation}
 |c(I)|\le 2^{|I|}\Bigl(\prod_{x\in I}p_x\Bigr)2^{\rho_*(I)},\qquad
 \rho_*(I)\le\rho_{\mathrm{rel}}(I).
 \label{p1:alg:compatible-bound}
\end{equation}
Equality holds in the second inequality when \(q(I)>0\).
\end{lemma}

\begin{proof}
The exact inclusion-exclusion formula is
\[
 c(I)=\sum_{A\subseteq I}(-1)^{|I|-|A|}q(A)
                       \prod_{x\in I\setminus A}p_x.
\]
A compatible term has absolute value
\(\bigl(\prod_{x\in I}p_x\bigr)2^{\rho_{\mathrm{rel}}(A)}\), and an incompatible term is
zero. There are at most \(2^{|I|}\) terms.
Nullity increases under inclusion, because adding the block
\(S_x\) increases the total rank by at most \(d_x\). This proves the other assertions.
\end{proof}

\subsection{Canonical direct components}

A partition \(\mathcal P\) of \(I\) is called direct if
the spaces \(V_C=\sum_{x\in C}S_x\), \(C\in\mathcal P\),
are in direct sum. Its finest such partition is the
canonical direct decomposition of \(I\).
We call a nonempty family \emph{connected} if this
partition has only one part.

\begin{lemma}[Existence, uniqueness, and factorization]
\label{p1:alg:components}
Every finite family \(I\) has a unique finest direct
partition \(\mathcal P(I)\). For it,
\begin{equation}
 q(I)=\prod_{C\in\mathcal P(I)}q(C),\qquad
 c(I)=\prod_{C\in\mathcal P(I)}c(C),\qquad
 \rho_{\mathrm{rel}}(I)=\sum_{C\in\mathcal P(I)}\rho_{\mathrm{rel}}(C).
 \label{p1:alg:factorization}
\end{equation}
\end{lemma}

\begin{proof}
If \(\mathcal P,\mathcal Q\) are direct partitions,
their nonempty intersections form a direct partition.
Indeed, inside each \(P\in\mathcal P\), the spaces
\(\sum_{x\in P\cap Q}S_x\) lie in the mutually direct
spaces \(V_Q\), and sum to \(V_P\). The direct sum over
\(P\) proves the claim. Intersecting the finitely many
direct partitions gives the unique finest one.

Characters indexed by mutually direct subspaces of
the finite prime-sign vector space are independent
uniform character families. Events, and centered
variables, depending on different \(V_C\) are therefore
independent. This proves both probability identities;
rank additivity proves the last one.
\end{proof}

Connectedness here is defined by all relative
relations. A minimal dependent family of original
blocks is primitive, but primitive block supports
alone do not determine the canonical components:
an additional relation can join a third block to
an already dependent pair. Exact pair compression
will retain this possibility.
